\SourcePage{260}
\SGAsection{4}{相对完美性}
\begin{SGAresult}{定义}{4.1}
设 $f:X\to Y$ 是\Term{scheme}{概形}的\Term{locally-finite-type}{局部有限型}\Term{morphism}{态射}，$E\in\operatorname{ob}D(X)$，$[a,b]$ 是 $\mathbf Z$ 中的区间。如果 $E$ 相对于 $f$ 是\Term{pseudo-coherent}{伪凝聚}的（1.2），且其\Term{flat-amplitude}{平坦振幅}包含于 $[a,b]$（3.1），就称 $E$ 相对于 $f$（无歧义时也说相对于 $Y$）的\Term{perfect-amplitude}{完美振幅}包含于 $[a,b]$，记作
\[
\operatorname{parf.amp}_f(E)\subset[a,b]
\qquad\text{（或 }\operatorname{parf.amp}_Y(E)\subset[a,b]\text{）。}
\]
如果 $\mathcal O_X$ 的\Term{perfect-amplitude}{完美振幅}包含于 $[-n,0]$，就称 $f$ 的\Term{perfect-amplitude}{完美振幅}包含于 $[-n,0]$，或称 $f$ 的\Term{perfect-dimension}{完美维数}不超过 $n$（$n\geq0$）。如果存在区间 $[a,b]$ 使 $\operatorname{parf.amp}_f(E)\subset[a,b]$，就称 $E$ 相对于 $f$ 具有\Term{finite}{有限}\Term{perfect-amplitude}{完美振幅}。如果 $E$ \Term{local-scheme}{局部}地相对于 $f$ 具有\Term{finite}{有限}\Term{perfect-amplitude}{完美振幅}，就称 $E$ 相对于 $f$ 是\Term{perfect-complex}{完美复形}。如果 $f$ \Term{local-scheme}{局部}地具有\Term{finite}{有限}\Term{perfect-amplitude}{完美振幅}，就称 $f$ 是\Term{perfect-morphism}{完美态射}。
\end{SGAresult}
\textbf{例 4.1.1。} 一个\Term{codimension}{余维}为 $d$ 的\Term{regular-immersion}{正则浸入}（第 VII 讲 1.4）的\Term{perfect-dimension}{完美维数}不超过 $d$（由 \Term{koszul-complex}{Koszul 复形}即可看出）。一个\Term{local-complete-intersection-morphism}{局部完全交态射}，即\Term{local-scheme}{局部}\Term{decomposable}{可分解}为一个\Term{regular-immersion}{正则浸入}与一个\Term{smooth}{光滑}\Term{morphism}{态射}的复合的\Term{morphism}{态射}，
\SourcePage{261}
是\Term{perfect-morphism}{完美态射}。

\textbf{记号 4.2。} 设 $f:X\to Y$ 是\Term{scheme}{概形}的\Term{locally-finite-type}{局部有限型}\Term{morphism}{态射}。用 $D(f)_{\mathrm{parf}}$、$D(X)_{f\text{-}\mathrm{parf}}$，或在无混淆之虞时用 $D(X)_{Y\text{-}\mathrm{parf}}$（相应地，用 $D(f)_{\underline{\mathrm{parf}}}$、$D(X)_{f\text{-}\underline{\mathrm{parf}}}$、$D(X)_{Y\text{-}\underline{\mathrm{parf}}}$），表示 $D(X)$ 中由相对于 $f$ 的\Term{perfect-complex}{完美复形}（相应地，具有\Term{finite}{有限}\Term{perfect-amplitude}{完美振幅}的\Term{complex}{复形}）组成的\Term{full-subcategory}{全子范畴}。由第 I 讲 5.3 和 1.4，$D(f)_{\mathrm{parf}}$ 与 $D(f)_{\underline{\mathrm{parf}}}$ 是 $D(X)$ 的\Term{triangulated-subcategory}{三角子范畴}。如果 $X$ \Term{quasi-compact}{拟紧}，则 $D(f)_{\mathrm{parf}}=D(f)_{\underline{\mathrm{parf}}}$。

\begin{SGAresult}{命题}{4.3}
（参见第 I 讲 5.8。）设 $f:X\to Y$ 是\Term{scheme}{概形}的\Term{locally-finite-type}{局部有限型}\Term{morphism}{态射}，$E\in\operatorname{ob}D(X)$，$[a,b]\subset\mathbf Z$ 为区间，$x\in X$，$y=f(x)$。以下条件\Term{equivalence}{等价}：
\begin{enumerate}
\item[(i)] 在 $x$ 的某个\Term{neighborhood}{邻域}上有 $\operatorname{parf.amp}_f(E)\subset[a,b]$；
\item[(ii)] 在 $x$ 的某个\Term{neighborhood}{邻域}上，$E$ 相对于 $f$ 是 $(a-1)$-\Term{pseudo-coherent}{伪凝聚}的，在 $[a,b]$ 以外\Term{acyclic}{无上同调}，且 $\operatorname{tor.amp}_{\mathcal O_{Y,y}}(E_x)\subset[a,b]$。
\end{enumerate}
\end{SGAresult}
\begin{proof}
(i)$\Rightarrow$(ii) 显然。证明 (ii)$\Rightarrow$(i)。问题在 $x$ 附近是\Term{local-scheme}{局部}的，故由定义立即归结到 $f$ \Term{smooth}{光滑}的情形。此时 $E$ 在 $x$ 附近是 $(a-1)$-\Term{pseudo-coherent}{伪凝聚}的，所以经\Term{localization}{局部化}，可设 $E$ 严格 $(a-1)$-\Term{pseudo-coherent}{伪凝聚}（第 I 讲 2.1），且在\Term{degree}{次数} $>b$ 时为零；再经\Term{truncation}{截断}，即将 $E$ 换成
\[
(0\longrightarrow E^a/B^a\longrightarrow E^{a+1}\longrightarrow\cdots\longrightarrow E^b\longrightarrow0),
\]
可设 $i>a$ 时 $E^i$ \Term{free-finite-type}{自由且有限型}，$E^a$ \Term{finitely-presented}{有限表示}，而 $i\notin[a,b]$ 时 $E^i=0$。由第 I 讲 5.1.1，$E_x^a$ 相对于 $\mathcal O_{Y,y}$ \Term{flat}{平坦}，故由 EGA IV 11.3.1，$E^a$ 在 $x$ 附近是 $f$-\Term{flat}{平坦}的。简记 $E^a=F$，须证 $F$ 在 $x$ 附近\Term{pseudo-coherent}{伪凝聚}。将 $X$ \Term{embedding}{嵌入} $Y$ 上的一个标准\Term{affine-space}{仿射空间}后，
\SourcePage{262}
可设 $X=Y[t_1,\ldots,t_r]$，其中 $Y$ \Term{affine}{仿射}，$f$ 为\Term{canonical-projection}{典范投影}。于是 $Y$ 是\Term{noetherian}{诺特}\Term{affine-scheme}{仿射概形} $Y_\alpha$ 的\Term{projective-limit}{投射极限}。记 $f_\alpha:X_\alpha=Y_\alpha[t_1,\ldots,t_r]\to Y_\alpha$ 为\Term{canonical-projection}{典范投影}。由 EGA IV 8.5.2 和 11.2.6，存在指标 $\alpha$ 以及\Term{finitely-presented}{有限表示}且 $f_\alpha$-\Term{flat}{平坦}的 $\mathcal O_{X_\alpha}$-\Term{module}{模} $F_\alpha$，使 $g_\alpha^*(F_\alpha)\simeq F$，其中 $g_\alpha:X\to X_\alpha$ 为\Term{canonical}{典范}\Term{morphism}{态射}。因 $X_\alpha$ \Term{noetherian}{诺特}，$F_\alpha$ \Term{coherent}{凝聚}，故相对于 $f_\alpha$ \Term{pseudo-coherent}{伪凝聚}。另一方面，$F_\alpha$ 是 $f_\alpha$-\Term{flat}{平坦}的，且 $f_\alpha$ \Term{flat}{平坦}，所以
\[
g_\alpha^*(F_\alpha)\simeq F_\alpha\mathbin{\overset{\mathbf L}{\otimes}}_{Y_\alpha}\mathcal O_Y.
\]
由 1.10 a)，$F$ 相对于 $Y$ \Term{pseudo-coherent}{伪凝聚}，因而\Term{pseudo-coherent}{伪凝聚}。证毕。
\end{proof}

\begin{SGAresult}{推论}{4.3.1}
设 $f:X\to Y$ 是\Term{locally-finite-type}{局部有限型}\Term{morphism}{态射}，$E$ 是 $\mathcal O_X$-\Term{module}{模}。$\operatorname{parf.amp}_f(E)\subset[0,0]$ 的充要条件是 $E$ 相对于 $f$ \Term{flat}{平坦}且\Term{finitely-presented}{有限表示}（1.2.1 c)）。特别地，$f$ 的\Term{perfect-dimension}{完美维数}为零（4.1）的充要条件是 $f$ \Term{flat}{平坦}且\Term{locally-finitely-presented}{局部有限表示}。
\end{SGAresult}
\begin{proof}在 4.3 的 (ii)$\Rightarrow$(i) 中取 $a=b=0$。
\end{proof}
\begin{SGAresult}{推论}{4.3.2}
设 $f:X\to Y$ 是\Term{locally-finite-type}{局部有限型}\Term{morphism}{态射}，$E\in\operatorname{ob}D(X)$ 是由相对于 $f$ \Term{flat}{平坦}且\Term{finitely-presented}{有限表示}的 $\mathcal O_X$-\Term{module}{模}组成的\Term{complex}{复形}（1.2.1 c)），$[a,b]\subset\mathbf Z$ 为区间。若 $i\notin[a,b]$ 时 $E^i=0$，则 $\operatorname{parf.amp}_f(E)\subset[a,b]$。若 $E$ 的\Term{degree}{次数}有上界，则 $E$ 相对于 $f$ \Term{pseudo-coherent}{伪凝聚}。
\end{SGAresult}
\begin{proof}由 4.3.1 以及第 I 讲 4.15、2.11 得出。
\end{proof}
\begin{SGAresult}{命题}{4.4}
设有\Term{scheme}{概形}的\Term{commutative-triangle}{交换三角形}
\[
\begin{tikzcd}
X\arrow[rr,hook,"i"]\arrow[dr,"f"']&&Y\arrow[dl,"g"]\\
&S&
\end{tikzcd}
\]
\SourcePage{263}
其中 $i$ 为\Term{closed-immersion}{闭浸入}，$g$ 为\Term{smooth}{光滑}\Term{morphism}{态射}，并设 $E\in\operatorname{ob}D(X)$。$E$ 相对于 $f$ 是\Term{perfect-complex}{完美复形}的充要条件是 $i_*(E)$ 为绝对意义下的\Term{perfect-complex}{完美复形}（第 I 讲 4.8）。更确切地，设 $g$ 的\Term{relative-dimension}{相对维数}不超过 $d$，$[a,b]\subset\mathbf Z$ 为区间，则
\begin{align*}
\text{(i)}\quad&\operatorname{parf.amp}(i_*(E))\subset[a,b]\Longrightarrow\operatorname{parf.amp}_f(E)\subset[a,b];\\
\text{(ii)}\quad&\operatorname{parf.amp}_f(E)\subset[a,b]\Longrightarrow\operatorname{parf.amp}(i_*(E))\subset[a-d,b].
\end{align*}
\end{SGAresult}
\begin{proof}这是 3.6 的直接推论。
\end{proof}
\begin{SGAresult}{命题}{4.5}
设
\[
\begin{tikzcd}
X\arrow[rr,"f"]\arrow[dr,"h"']&&Y\arrow[dl,"g"]\\
&S&
\end{tikzcd}
\]
为\Term{scheme}{概形}的\Term{locally-finite-type}{局部有限型}\Term{morphism}{态射}的\Term{commutative-triangle}{交换三角形}，并设 $E\in\operatorname{ob}D^-(X)$、$F\in\operatorname{ob}D^-(Y)$，$[a,b]$、$[c,d]$ 为 $\mathbf Z$ 中的区间。则
\[
\operatorname{parf.amp}_f(E)\subset[a,b],\quad\operatorname{parf.amp}_g(F)\subset[c,d]
\Longrightarrow
\operatorname{parf.amp}_h(E\mathbin{\overset{\mathbf L}{\otimes}}_YF)\subset[a+c,b+d].
\]
因此，若 $E$ 相对于 $f$ 是\Term{perfect-complex}{完美复形}，$F$ 相对于 $g$ 是\Term{perfect-complex}{完美复形}，则 $E\mathbin{\overset{\mathbf L}{\otimes}}_YF$ 相对于 $h$ 是\Term{perfect-complex}{完美复形}。
\end{SGAresult}
\begin{proof}结合 1.11 与 3.4。
\end{proof}
\begin{SGAresult}{推论}{4.5.1}
在 4.5 的假设下：
\begin{enumerate}
\item[a)] 若 $g$ 的\Term{perfect-dimension}{完美维数}不超过 $n$，则
\[
\operatorname{parf.amp}_f(E)\subset[a,b]\Longrightarrow\operatorname{parf.amp}_h(E)\subset[a-n,b];
\]
因此，若 $g$ 是\Term{perfect-morphism}{完美态射}，则 $E$ 相对于 $f$ 为\Term{perfect-complex}{完美复形}蕴含 $E$ 相对于 $h$ 为\Term{perfect-complex}{完美复形}。
\item[b)] 若 $f$ 的\Term{perfect-dimension}{完美维数}不超过 $m$，则
\[
\operatorname{parf.amp}_g(F)\subset[c,d]\Longrightarrow\operatorname{parf.amp}_h(\mathbf Lf^*(F))\subset[c-m,d];
\]
\SourcePage{264}
因此，若 $f$ 是\Term{perfect-morphism}{完美态射}，则 $F$ 相对于 $g$ 为\Term{perfect-complex}{完美复形}蕴含 $\mathbf Lf^*(F)$ 相对于 $h$ 为\Term{perfect-complex}{完美复形}。
\end{enumerate}
\end{SGAresult}
\textbf{注记 4.5.2。} 沿用上述记号，即使 $g$ 是\Term{perfect-morphism}{完美态射}，$E$ 相对于 $h$ 为\Term{perfect-complex}{完美复形}也未必蕴含它相对于 $f$ 为\Term{perfect-complex}{完美复形}，域上\Term{scheme}{概形}的\Term{singular-point}{奇点}这个反复出现的例子就说明了这一点。有如下判据。

\begin{SGAresult}{命题}{4.6}
（逐\Term{fiber}{纤维}\Term{perfect}{完美}性判据。）\footnote{这是“\Term{fibrewise-flatness-criterion}{逐纤维平坦性判据}”（EGA IV 11.3.10）的类似物。}设
\[
\begin{tikzcd}
X\arrow[rr,"f"]\arrow[dr,"h"']&&Y\arrow[dl,"g"]\\
&S&
\end{tikzcd}
\]
为\Term{scheme}{概形}\Term{morphism}{态射}的\Term{commutative-triangle}{交换三角形}。假设 $f,h$ \Term{locally-finite-type}{局部有限型}，$g$ \Term{locally-finitely-presented}{局部有限表示}且\Term{flat}{平坦}。设 $x\in X$，$y=f(x)$，$s=h(x)=g(y)$。又设 $E\in\operatorname{ob}D^-(X)$，$[a,b]\subset\mathbf Z$ 为区间。记 $f_s:X_s\to Y_s$ 为 $f$ 在 $s$ 处的\Term{fiber}{纤维}，即经\Term{base-change}{基变换} $\operatorname{Spec}(k(s))\to S$ 由 $f$ 得到的\Term{morphism}{态射}。以下条件\Term{equivalence}{等价}：
\begin{enumerate}
\item[(i)] 在 $x$ 附近 $\operatorname{parf.amp}_h(E)\subset[a,b]$，且
\[
\operatorname{tor.amp}_{\mathcal O_{Y_s,y}}
\bigl(E_x\mathbin{\overset{\mathbf L}{\otimes}}_{\mathcal O_{Y,y}}\mathcal O_{Y_s,y}\bigr)\subset[a,b];
\]
\item[(ii)] 在 $x$ 附近 $\operatorname{parf.amp}_f(E)\subset[a,b]$。
\end{enumerate}
\end{SGAresult}
\begin{proof}
(ii)$\Rightarrow$(i)。已知 $g$ 的\Term{perfect-dimension}{完美维数}为零（4.3.1）。故由 4.5.1 a)，条件 (ii) 蕴含在 $x$ 附近 $\operatorname{parf.amp}_h(E)\subset[a,b]$。由 4.3，它还蕴含 $\operatorname{tor.amp}_{\mathcal O_{Y,y}}(E_x)\subset[a,b]$，于是由第 I 讲 5.1 (i)，
\[
\operatorname{tor.amp}_{\mathcal O_{Y_s,y}}
\bigl(E_x\mathbin{\overset{\mathbf L}{\otimes}}_{\mathcal O_{Y,y}}\mathcal O_{Y_s,y}\bigr)\subset[a,b],
\]
\SourcePage{265}
从而 (ii)$\Rightarrow$(i) 得证。

(i)$\Rightarrow$(ii)。问题在 $x$ 附近是\Term{local-scheme}{局部}的，可设 $f$ \Term{resolution}{分解}为 $X\xrightarrow{i}X'\xrightarrow{f'}Y$，其中 $i$ 为\Term{closed-immersion}{闭浸入}，$f'$ 为\Term{smooth}{光滑}\Term{morphism}{态射}。将 $E$ 换成 $i_*(E)$、$f$ 换成 $f'$（并将 $h$ 换成 $gf'$）显然不改变条件 (i)、(ii)，因此可设 $f$ \Term{smooth}{光滑}。因 $g$ 是\Term{perfect-morphism}{完美态射}（4.3.1），$h$ 也是\Term{perfect-morphism}{完美态射}，所以由 1.12，$E$ \Term{pseudo-coherent}{伪凝聚}。经\Term{localization}{局部化}，由第 II 讲 2.2.10 b) 和第 I 讲 2.10 a)，可设 $E$ 严格\Term{pseudo-coherent}{伪凝聚}（第 I 讲 2.1），在\Term{degree}{次数} $>b$ 时为零，且\Term{truncation}{截断}\Term{arrow}{箭头}
\[
t:E\longrightarrow(0\longrightarrow E^a/B^a\longrightarrow E^{a+1}\longrightarrow\cdots\longrightarrow E^b\longrightarrow0)
\]
为\Term{quasi-isomorphism}{拟同构}（事实上，假设 (i) 蕴含 $E$ 在 $s$ 附近在 $[a,b]$ 以外\Term{acyclic}{无上同调}）。须证 $E^a/B^a$ 在 $x$ 附近是 $f$-\Term{flat}{平坦}的。记 $i:X_s\to X$ 为\Term{canonical}{典范}\Term{arrow}{箭头}。有
\[
E_x\mathbin{\overset{\mathbf L}{\otimes}}_{\mathcal O_{Y,y}}\mathcal O_{Y_s,y}\simeq i^*(E)_x.
\]
因 $h$ \Term{flat}{平坦}，由 3.5.2，假设 (i) 蕴含在 $x$ 附近 $\operatorname{tor.amp}_{h_s}(i^*(E))\subset[a,b]$。所以经\Term{localization}{局部化}，可设\Term{truncation}{截断}\Term{arrow}{箭头}
\[
i^*(t):i^*(E)\longrightarrow
(0\longrightarrow i^*(E^a)/B^a(i^*(E))\longrightarrow\cdots\longrightarrow i^*(E^b)\longrightarrow0)
\]
为\Term{quasi-isomorphism}{拟同构}。于是由第 I 讲 5.1.1，假设 $\operatorname{tor.amp}_{\mathcal O_{Y_s,y}}i^*(E)_x\subset[a,b]$ 蕴含 $i^*(E^a/B^a)_x$ 是\Term{flat}{平坦} $\mathcal O_{Y_s,y}$-\Term{module}{模}。

另一方面，因\Term{arrow}{箭头} $(*)$ 为\Term{quasi-isomorphism}{拟同构}，仍由第 I 讲 5.1.1，假设 $\operatorname{tor.amp}_h(E)\subset[a,b]$ 蕴含
\SourcePage{266}
$E^a/B^a$ 是 $h$-\Term{flat}{平坦}的。由\Term{fibrewise-flatness-criterion}{逐纤维平坦性判据}（EGA IV 11.3.10），得出 $E^a/B^a$ 在 $x$ 附近是 $f$-\Term{flat}{平坦}的。证毕。
\end{proof}
\begin{SGAresult}{推论}{4.6.1}
设 $f:X\to S$ 是\Term{locally-finitely-presented}{局部有限表示}且\Term{flat}{平坦}的\Term{morphism}{态射}，$x\in X$，$s=f(x)$，$E\in\operatorname{ob}D^-(X)$，$[a,b]\subset\mathbf Z$ 为区间。记 $i_s:X_s=X\times_S\operatorname{Spec}(k(s))\to X$ 为\Term{canonical-projection}{典范投影}，并置 $E_{X_s}=\mathbf Li_s^*(E)$。则
\[
\left\{
\begin{gathered}
\operatorname{parf.amp}_f(E)\subset[a,b]\quad\text{在 }x\text{ 附近},\\
\operatorname{tor.amp}(E_{X_s})_x\subset[a,b]
\end{gathered}
\right.
\quad\Longleftrightarrow\quad
\operatorname{parf.amp}(E)\subset[a,b]\quad\text{在 }x\text{ 附近}.
\]
因此，$E$ 为\Term{perfect-complex}{完美复形}，当且仅当 $E$ 相对于 $f$ 为\Term{perfect-complex}{完美复形}，且对每个 $s\in S$，$E_{X_s}$ 均为\Term{perfect-complex}{完美复形}。
\end{SGAresult}
\begin{proof}在 4.6 中取 $X=Y$、$f=\operatorname{Id}_X$。
\end{proof}
\textbf{注记 4.6.2。} 由 4.6.1 和第 I 讲 5.10，若 $f$ \Term{locally-finitely-presented}{局部有限表示}、\Term{flat}{平坦}，且各\Term{fiber}{纤维}均为\Term{regular-scheme}{正则概形}，则 $E$ 相对于 $f$ 为\Term{perfect-complex}{完美复形}\Term{equivalence}{等价}于 $E$ 为\Term{perfect-complex}{完美复形}（对于 $f$ \Term{smooth}{光滑}的情形，此\Term{equivalence-relation}{等价关系}已在 4.4 中得到）。

\begin{SGAresult}{命题}{4.7}
设 $S$ 为\Term{scheme}{概形}，对 $i=1,2$，设 $f_i:X_i\to Y_i$ 为\Term{locally-finite-type}{局部有限型} $S$-\Term{morphism}{态射}，$E_i\in\operatorname{ob}D(X_i)$，$[a_i,b_i]\subset\mathbf Z$ 为区间。假设 $X_1,X_2$ 以及 $Y_1,Y_2$ 在 $S$ 上均\Term{tor-independent}{Tor 无关}（1.5）。置 $a=a_1+a_2$、$b=b_1+b_2$、$f=f_1\times_Sf_2$。则：
\begin{enumerate}
\item[a)] 若 $\operatorname{parf.amp}_{f_1}(E_1)\subset[a_1,b_1]$，且 $E_2$ 相对于 $f_2$ \Term{pseudo-coherent}{伪凝聚}并在 $[a_2,b_2]$ 以外\Term{acyclic}{无上同调}，则 $E_1\mathbin{\overset{\mathbf L}{\otimes}}_SE_2$（参见 1.6）相对于 $f$ \Term{pseudo-coherent}{伪凝聚}，并在 $[a,b]$ 以外\Term{acyclic}{无上同调}。
\SourcePage{267}
\item[b)] 对 $i=1,2$ 有 $\operatorname{parf.amp}_{f_i}(E_i)\subset[a_i,b_i]$，则
\[
\operatorname{parf.amp}_f(E_1\mathbin{\overset{\mathbf L}{\otimes}}_SE_2)\subset[a,b].
\]
\end{enumerate}
\end{SGAresult}
\begin{proof}结合 1.9 和 3.5。
\end{proof}
\begin{SGAresult}{推论}{4.7.1}
在 4.6 的假设下，若对 $i=1,2$，$f_i$ 的\Term{perfect-dimension}{完美维数}不超过 $n_i$，则 $f$ 的\Term{perfect-dimension}{完美维数}不超过 $n_1+n_2$。因此，若 $f_1,f_2$ 都是\Term{perfect-morphism}{完美态射}，则 $f$ 也是。
\end{SGAresult}
\begin{SGAresult}{推论}{4.7.2}
设有\Term{scheme}{概形}的\Term{cartesian-square}{笛卡尔方块}
\[
\begin{tikzcd}
X\arrow[d,"f"']&X'\arrow[l,"g"']\arrow[d,"f'"]\\
Y&Y'\arrow[l,"h"]
\end{tikzcd}
\]
其中 $f$ \Term{locally-finite-type}{局部有限型}，$X$ 和 $Y'$ 在 $Y$ 上\Term{tor-independent}{Tor 无关}。设 $E\in\operatorname{ob}D(X)$，$[a,b]\subset\mathbf Z$ 为区间。则
\[
\operatorname{parf.amp}_f(E)\subset[a,b]\Longrightarrow
\operatorname{parf.amp}_{f'}(\mathbf Lg^*(E))\subset[a,b].
\]
因此，\Term{functor}{函子} $\mathbf Lg^*:D^-(X)\to D^-(X')$ 将 $D(f)_{\mathrm{parf}}$（相应地 $D(f)_{\underline{\mathrm{parf}}}$）映入 $D(f')_{\mathrm{parf}}$（相应地 $D(f')_{\underline{\mathrm{parf}}}$）（4.2）。
\end{SGAresult}
\begin{SGAresult}{命题}{4.8}
设 $S$ 为\Term{scheme}{概形}，$f:X\to Y$ 是\Term{locally-finite-type}{局部有限型} $S$-\Term{scheme}{概形}之间的\Term{proper}{固有} $S$-\Term{morphism}{态射}，$E\in\operatorname{ob}D^b(X)$。假定 2.1 对 $f$ 成立（所以当 $S$ \Term{locally-noetherian}{局部诺特}或 $f$ \Term{projective}{射影}时总成立），则：
\begin{enumerate}
\item[a)] 若 $E$ 相对于 $S$ 为\Term{perfect-complex}{完美复形}，则 $\mathbf Rf_*(E)$ 也如此。
\item[b)] 设 $S$ \Term{quasi-compact}{拟紧}，$X,Y$ 在 $S$ 上\Term{finite-type}{有限型}，$[a,b]\subset\mathbf Z$ 为区间，$N$ 为使任意\Term{quasi-coherent}{拟凝聚} $\mathcal O_X$-\Term{module}{模} $M$ 在 $i>N$ 时都有 $R^if_*(M)=0$ 的整数
\SourcePage{268}
（由 EGA III 1.4.12 和 IV 1.7.21，这样的整数存在），则
\[
\operatorname{parf.amp}_S(E)\subset[a,b]\Longrightarrow
\operatorname{parf.amp}_S(\mathbf Rf_*(E))\subset[a,b+N].
\]
因此，\Term{functor}{函子} $\mathbf Rf_*$ 将 $D(X)_{S\text{-}\underline{\mathrm{parf}}}$ 映入 $D(Y)_{S\text{-}\underline{\mathrm{parf}}}$。\footnote{事实上，由于 $X,Y$ \Term{quasi-compact}{拟紧}，有 $D(X)_{S\text{-}\underline{\mathrm{parf}}}=D(X)_{S\text{-}\mathrm{parf}}$，对 $Y$ 也一样（参见 4.2）。}
\end{enumerate}
\end{SGAresult}
\begin{proof}断言 a) 由 b) 得出；b) 则由结合 2.1 与 3.7.1 得到。
\end{proof}
\begin{SGAresult}{推论}{4.8.1}
设 $f:X\to Y$ 是\Term{proper}{固有}且\Term{perfect-morphism}{完美}的\Term{morphism}{态射}，$E\in\operatorname{ob}D^+(X)$。假定 2.1 对 $f$ 成立，则若 $E$ 是\Term{perfect-complex}{完美复形}，$\mathbf Rf_*(E)$ 也如此；此外，若 $Y$ \Term{quasi-compact}{拟紧}，则\Term{functor}{函子} $\mathbf Rf_*$ 将 $D(X)_{\underline{\mathrm{parf}}}$ 映入 $D(Y)_{\underline{\mathrm{parf}}}$。\footnote{有 $D(X)_{\underline{\mathrm{parf}}}=D(X)_{\mathrm{parf}}$，对 $Y$ 也一样。}
\end{SGAresult}
\begin{proof}事实上，由 4.5.1 a)，若 $E$ 是\Term{perfect-complex}{完美复形}，则它相对于 $f$ 也是\Term{perfect-complex}{完美复形}。
\end{proof}
\begin{SGAresult}{命题}{4.9}
\footnote{本结果及其推论在本讨论班后文中不再使用，但它们在\Term{coherent}{凝聚}\Term{sheaf}{层}的 Lefschetz--Verdier 公式中起重要\Term{action}{作用}，参见 SGA 5 第 III 讲。}设 $S$ 是\Term{noetherian}{诺特}\Term{scheme}{概形}，$X,Y$ 是 $S$ 上的\Term{finite-type}{有限型}\Term{scheme}{概形}，$f:X\to Y$ 是\Term{compactifiable}{可紧化（compactifiable）}的 $S$-\Term{morphism}{态射}，即一个\Term{open-immersion}{开浸入}与一个\Term{proper}{固有}\Term{morphism}{态射}的复合。设 $E\in\operatorname{ob}D(X)$、$F\in\operatorname{ob}D^+(Y)$。若 $E$ 相对于 $f$ 为\Term{perfect-complex}{完美复形}，$F$ 相对于 $S$ 为\Term{perfect-complex}{完美复形}，则 $\mathbf R\mathcal Hom(E,f^!(F))$ 相对于 $S$ 为\Term{perfect-complex}{完美复形}。\footnote{\Term{functor}{函子} $f^!$ 的定义见 [H] 附录第 4 节。}
\end{SGAresult}
\begin{proof}
问题在 $X$ 上是\Term{local-scheme}{局部}的，可设 $f$ \Term{resolution}{分解}为 $X\xrightarrow{i}X'\xrightarrow{f'}Y$，其中 $i$ 是\Term{closed-immersion}{闭浸入}，$f'$ 是\Term{smooth}{光滑}\Term{morphism}{态射}。若 $E$ 相对于 $f$ 为\Term{perfect-complex}{完美复形}，则它相对于 $i$ 也为\Term{perfect-complex}{完美复形}（4.4）；若 $F$ 相对于 $S$ 为\Term{perfect-complex}{完美复形}，则由 4.5.1 b)，$f'^!(F)$ 相对于 $S$ 为\Term{perfect-complex}{完美复形}，因为 $f'^!(F)$ \Term{local-scheme}{局部}地\Term{isomorphism}{同构}于 $f'^*(F)$ 的一个\Term{translation}{平移}。又 $f^!\simeq i^!f'^!$，故归结到 $f$ 为\Term{closed-immersion}{闭浸入}的情形。于是设 $E$ 相对于 $f$ 为\Term{perfect-complex}{完美复形}，即 $f_*(E)$ 为\Term{perfect-complex}{完美复形}，且 $F$ 相对于 $S$ 为\Term{perfect-complex}{完美复形}。须证 $f_*\mathbf R\mathcal Hom(E,f^!F)$ 相对于 $S$ 为
\SourcePage{269}
\Term{perfect-complex}{完美复形}。由\Term{duality-theorem}{对偶定理}（[H] III 6.7），
\[
f_*\mathbf R\mathcal Hom(E,f^!F)\simeq\mathbf R\mathcal Hom(f_*(E),F).
\]
由于 $f_*(E)$ 是\Term{perfect-complex}{完美复形}，由第 I 讲 7.7，
\[
\mathbf R\mathcal Hom(f_*(E),F)\simeq(f_*(E))^\vee\mathbin{\overset{\mathbf L}{\otimes}}F,
\]
再用 4.5（其中 $X=Y$ 且 $f=\operatorname{Id}_X$）即得结论。
\end{proof}
\begin{SGAresult}{推论}{4.9.1}
在 4.8 的假设下，若 $f$ 是\Term{perfect-morphism}{完美态射}，则 $f^!$ 将 $D(Y)_{S\text{-}\mathrm{parf}}$ 映入 $D(X)_{S\text{-}\mathrm{parf}}$。
\end{SGAresult}
\begin{proof}取 $E=\mathcal O_X$。
\end{proof}
\begin{SGAresult}{推论}{4.9.2}
设 $S$ 为\Term{noetherian}{诺特}\Term{scheme}{概形}，$f:X\to S$ 为\Term{compactifiable}{可紧化}\Term{morphism}{态射}。置 $f^!(\mathcal O_S)=K_X$，$\mathbf R\mathcal Hom(\,\cdot\,,K_X)=\underline D_X$。对 $E\in\operatorname{ob}D(f)_{\mathrm{parf}}$，有 $\underline D_X(E)\in\operatorname{ob}D(f)_{\mathrm{parf}}$，且\Term{canonical}{典范}\Term{arrow}{箭头}（第 I 讲 7.2）
\[
E\longrightarrow\underline D_X\underline D_X(E)
\]
为\Term{isomorphism}{同构}。
\end{SGAresult}
\begin{proof}
第一个断言是在 4.8 中取 $Y=S$、$F=\mathcal O_S$ 的结果。对第二个断言，容易归结到 $f$ 为\Term{closed-immersion}{闭浸入}的情形。此时只须证明对\Term{canonical}{典范}\Term{arrow}{箭头}\Term{action}{作用} $f_*$ 后所得的\Term{arrow}{箭头}为\Term{isomorphism}{同构}。由整体\Term{duality}{对偶}（[H] III 6.7），有 $f_*\underline D_X\simeq\underline D_Sf_*$，故归结到证明：当 $E$ 相对于 $f$ 为\Term{perfect-complex}{完美复形}，即 $f_*(E)$ 为\Term{perfect-complex}{完美复形}时，\Term{canonical}{典范}\Term{arrow}{箭头}
\[
f_*(E)\longrightarrow\underline D_S\underline D_Sf_*(E)
\]
为\Term{isomorphism}{同构}。这由\Term{perfect-complex}{完美复形}通常的\Term{biduality}{双对偶性}（第 I 讲 7.2）得出。
\end{proof}

\SGAsection{5}{应用：交换定理与半连续性定理}
\footnote{本节的结果在本讨论班后文中不再使用。}
\begin{SGAresult}{\Term{lemma}{引理}}{5.1}
设 $(X,\mathcal O_X)$ 为（交换）\Term{ringed-topos}{赋环拓扑斯}，记 $\operatorname{Mod}(X)$ 为 $\mathcal O_X$-\Term{module}{模}的\Term{category}{范畴}，$D(X)=D(\operatorname{Mod}(X))$。设 $E\in\operatorname{ob}D^-(X)$、$p\in\mathbf Z$。考虑以下条件：
\SourcePage{270}
\begin{enumerate}
\item[a)] 从 $\operatorname{Mod}(X)$ 到 $\operatorname{Mod}(X)$ 的\Term{functor}{函子} $M\mapsto H^p(E\mathbin{\overset{\mathbf L}{\otimes}}M)$ \Term{right-exact}{右正合}。
\item[a')] 对一切 $M\in\operatorname{ob}\operatorname{Mod}(X)$，\Term{canonical}{典范}\Term{arrow}{箭头}
\[
H^p(E)\otimes M\longrightarrow H^p(E\mathbin{\overset{\mathbf L}{\otimes}}M)
\]
为\Term{isomorphism}{同构}。
\item[b)] 从 $\operatorname{Mod}(X)$ 到 $\operatorname{Mod}(X)$ 的\Term{functor}{函子} $M\mapsto H^{p+1}(E\mathbin{\overset{\mathbf L}{\otimes}}M)$ \Term{left-exact}{左正合}。
\item[b')] 对一切 $M\in\operatorname{ob}\operatorname{Mod}(X)$，下面定义的\Term{canonical}{典范}\Term{arrow}{箭头}
\[
H^{p+1}(E\mathbin{\overset{\mathbf L}{\otimes}}M)\longrightarrow
\mathcal Hom(H^{-p-1}(E^\vee),M)\qquad(*)
\]
为\Term{isomorphism}{同构}。
\end{enumerate}
则有 $\text{a)}\Longleftrightarrow\text{a')}\Longleftrightarrow\text{b)}\Longleftarrow\text{b')}$。此外，若 $E$ \Term{pseudo-coherent}{伪凝聚}，则 a)、a')、b)、b') \Term{equivalence}{等价}。
\end{SGAresult}
\textbf{$(*)$ 的定义。} \Term{canonical}{典范}\Term{arrow}{箭头}（第 I 讲 7.5）$E^\vee\mathbin{\overset{\mathbf L}{\otimes}}E\to\mathcal O_X$，其中 $E^\vee=\mathbf R\mathcal Hom(E,\mathcal O_X)$，与 $M\in\operatorname{ob}\operatorname{Mod}(X)$ 作\Term{tensor-product}{张量积}，给出\Term{arrow}{箭头} $E^\vee\mathbin{\overset{\mathbf L}{\otimes}}E\mathbin{\overset{\mathbf L}{\otimes}}M\to M$，从而由 Cartan 所钟爱的\Term{adjunction-formula}{伴随公式}（第 I 讲 7.4），得到\Term{canonical}{典范}\Term{arrow}{箭头}
\[
E\mathbin{\overset{\mathbf L}{\otimes}}M\longrightarrow\mathbf R\mathcal Hom(E^\vee,M).\qquad(**)
\]
将 $(**)$ 在 $H^{p+1}$ 上诱导的\Term{arrow}{箭头}与\Term{canonical}{典范}\Term{arrow}{箭头}
\[
\mathcal Ext^{p+1}(E^\vee,M)\longrightarrow\mathcal Hom(H^{-p-1}(E^\vee),M)\qquad(***)
\]
复合，就得到 $(*)$。
\begin{proof}
蕴含 $\text{a')}\Rightarrow\text{a)}\Longleftrightarrow\text{b)}\Longleftarrow\text{b')}$ 显然。证明 a)$\Rightarrow$a')。设 $M$ 为 $\mathcal O_X$-\Term{module}{模}。存在\Term{exact-sequence}{正合列} $M^1\to M^0\to M\to0$，其中 $M^0,M^1$ 是形如 $\mathcal O_{U,X}$ 的\Term{sheaf}{层}的\Term{direct-sum}{直和}，
\SourcePage{271}
这里 $U$ 为 $X$ 的\Term{object}{对象}。利用\Term{functor}{函子} $H^p(E)\otimes\,\cdot\,$ 与 $H^p(E\mathbin{\overset{\mathbf L}{\otimes}}\,\cdot\,)$ 的\Term{right-exact}{右正合}性及其与\Term{direct-sum}{直和}的\Term{commutativity}{交换性}，只须对 $M=\mathcal O_{U,X}$、$U\in\operatorname{ob}X$ 证明 a') 中的\Term{arrow}{箭头}为\Term{isomorphism}{同构}。而 $H^p(E)\otimes\mathcal O_{U,X}$（相应地 $H^p(E\mathbin{\overset{\mathbf L}{\otimes}}\mathcal O_{U,X})$）\Term{canonical}{典范}地等同于 $i_!(H^p(E)|_U)$（相应地 $H^p(i_!(E|_U))$）（SGA 4 第 IV 讲 2.13），其中 $i_!$ 为 $U$ 外的\Term{extension-by-zero}{零延拓}；因 $i_!$ \Term{exact}{正合}，结论即得。

设 $E$ \Term{pseudo-coherent}{伪凝聚}，证明 b)$\Rightarrow$b')。显然只须对\Term{injective-object}{内射} $M$ 证明 $(*)$ 为\Term{isomorphism}{同构}。因 $M$ \Term{injective-object}{内射}时 $(***)$ 为\Term{isomorphism}{同构}，所要的结论由以下断言得出：
\begin{SGAresult}{断言}{5.1.1}
若 $E$ \Term{pseudo-coherent}{伪凝聚}、$M$ \Term{injective-object}{内射}，则\Term{arrow}{箭头} $(**)$ 为\Term{isomorphism}{同构}。
\end{SGAresult}
\textbf{5.1.1 的证明。} 记 $T$（相应地 $T'$）为从 $D^-(X)$ 到 $D(X)$ 的\Term{functor}{函子} $F\mapsto F\mathbin{\overset{\mathbf L}{\otimes}}M$（相应地 $\mathbf R\mathcal Hom(F^\vee,M)$）。设 $n\in\mathbf Z$；若 $i>n$ 时 $H^i(F)=0$，则 $i>n$ 时 $H^iT(F)=0$，$i<-n$ 时 $H^i(F^\vee)=0$，且在 $M$ \Term{injective-object}{内射}时，$i>n$ 时 $H^iT'(F)=0$。又因为\Term{arrow}{箭头} $(**)$ 与到\Term{open-object}{开对象}的\Term{restriction}{限制}交换，且\Term{injective-object}{内射对象}\Term{restriction}{限制}到\Term{open-object}{开对象}仍\Term{injective-object}{内射}，所以经\Term{localization}{局部化}和\Term{truncation}{截断}，可设 $E$ \Term{strictly-perfect}{严格完美}（2.1）。再经逐步\Term{reduced}{约化}，可归结到 $E=\mathcal O_X$，结论即得。
\end{proof}
\textbf{注记 5.2。} (i) 条件 a)（相应地 a')、b)、b')）是\Term{local-nature}{局部性质}：若 $(X_\alpha)$ 是 $X$ 的\Term{final-object}{终对象}的\Term{open-covering}{开覆盖}，且对每个 $\alpha$，$X_\alpha$ 和 $E|_{X_\alpha}$ 满足条件 a)（相应地其余条件），则 a)（相应地其余条件）成立。

(ii) 若 $x$ 为 $X$ 的一个点，记 a)$_x$（相应地 a')$_x$ 等）为关于 $E_x$ 及 $\operatorname{Mod}(\mathcal O_{X,x})$ 的相应条件。
\SourcePage{272}
若 $X$ 有\Term{enough-points}{足够多的点}，则条件 a)（相应地 a')、b)）\Term{equivalence}{等价}于“对 $X$ 的每个点 $x$，a)$_x$（相应地 a')$_x$、b)$_x$）成立”；若 $E$ \Term{pseudo-coherent}{伪凝聚}，还\Term{equivalence}{等价}于“对 $X$ 的每个点 $x$，b')$_x$ 成立”。此外，条件 a') \Term{equivalence}{等价}于将其中“对一切 $M\in\operatorname{ob}\operatorname{Mod}(X)$”换成“对一切形如 $\mathcal O_{U,X}^{p}\to\mathcal O_{U,X}^{q}$ 的\Term{homomorphism}{同态}的\Term{cokernel}{余核}的 $\mathcal O_X$-\Term{module}{模}，其中 $U$ 为 $X$ 的\Term{object}{对象}”所得的相应条件。

(iii) EGA III 7.4.2 的判据可逐字推广到这里：若 $E$ 的各\Term{set-component}{分量}\Term{flat}{平坦}，则条件 a) 仅表示 $\operatorname{Coker}(E^p\to E^{p+1})$ \Term{flat}{平坦}；所以若进一步假定 $E^p,E^{p+1}$ \Term{finitely-presented}{有限表示}，此\Term{cokernel}{余核}\Term{local-scheme}{局部}地是\Term{free-finite-type}{有限型自由}\Term{module}{模}的\Term{direct-summand}{直和因子}（第 I 讲 5.8.3）。特别地，由此得出：若 $E$ 是\Term{perfect-complex}{完美复形}，则 $E$ 对整数 $p$ 满足条件 a)，当且仅当 $E^\vee$ 对整数 $-p-1$ 满足条件 a)。

(iv) 设 $X$ 为\Term{quasi-separated}{拟分离}\Term{scheme}{概形}，$\mathcal O_X$ 为其\Term{structure-sheaf}{结构层}，$E\in\operatorname{ob}D^-(X)$ 是\Term{cohomology}{上同调}\Term{local-scheme}{局部}\Term{bounded}{有界}且\Term{quasi-coherent}{拟凝聚}的\Term{complex}{复形}。记 a$_{\mathrm{qcoh}}$)（相应地 a'$_{\mathrm{qcoh}}$) 等）为将 a)（相应地 a') 等）中的 $\operatorname{Mod}(X)$ 换成\Term{quasi-coherent}{拟凝聚} $\mathcal O_X$-\Term{module}{模}的\Term{category}{范畴} $\operatorname{Qcoh}(X)$ 后的相应条件。则
\[
\text{a)}\Longleftrightarrow\text{a}_{\mathrm{qcoh}}\text{)}
\Longleftrightarrow\text{a}'_{\mathrm{qcoh}}\text{)}
\Longleftrightarrow\text{b}_{\mathrm{qcoh}}\text{)}
\Longleftarrow\text{b}'_{\mathrm{qcoh}}\text{)},
\]
而若 $E$ \Term{pseudo-coherent}{伪凝聚}，则这些条件\Term{equivalence}{等价}。

事实上，因 $X$ \Term{quasi-separated}{拟分离}，$X$ 的任意\Term{affine-open}{仿射开集}上的\Term{quasi-coherent}{拟凝聚}\Term{sheaf}{层}均由 $X$ 上的一个\Term{quasi-coherent}{拟凝聚}\Term{sheaf}{层}诱导（EGA I 9.4.2 和 IV 1.2.7）。因此，条件 a$_{\mathrm{qcoh}}$)（相应地 a'$_{\mathrm{qcoh}}$)、b$_{\mathrm{qcoh}}$)、b'$_{\mathrm{qcoh}}$)）成立，当且仅当它在每个\Term{affine-open}{仿射开集}上成立。所以为证明上述蕴含，可设 $X$ \Term{affine}{仿射}。此时由第 II 讲 3.5，$E$ 在 $D(X)$ 中\Term{isomorphism}{同构}于一个\Term{quasi-coherent}{拟凝聚}\Term{module}{模}的\Term{complex}{复形}，而蕴含
\[
\text{a}_{\mathrm{qcoh}}\text{)}\Longleftrightarrow
\text{a}'_{\mathrm{qcoh}}\text{)}\Longleftrightarrow
\text{b}_{\mathrm{qcoh}}\text{)}\Longleftarrow
\text{b}'_{\mathrm{qcoh}}\text{)}
\]
\SourcePage{273}
由对\Term{ringed}{赋环}为 $\Gamma(X,\mathcal O_X)$ 的“单点”\Term{topos}{拓扑斯}应用 5.1 得出；由第 II 讲 2.2.10 a)，在 $E$ \Term{pseudo-coherent}{伪凝聚}时，b$_{\mathrm{qcoh}}$)$\Rightarrow$b'$_{\mathrm{qcoh}}$) 也同样得出。最后，\Term{equivalence-relation}{等价关系} a)$\Longleftrightarrow$a$_{\mathrm{qcoh}}$)，或\Term{equivalence}{等价}地 a')$\Longleftrightarrow$a'$_{\mathrm{qcoh}}$)，由 (ii) 的最后一条注记得出（并用到 $X$ \Term{quasi-separated}{拟分离}时\Term{quasi-coherent}{拟凝聚}\Term{sheaf}{层}“可延拓”）。

\begin{SGAresult}{命题}{5.3}
（\Term{commutativity}{交换性}质；EGA III 7.7.5 II。）设 $f:X\to Y$ 为\Term{quasi-compact}{拟紧}且\Term{quasi-separated}{拟分离}\Term{scheme}{概形}之间的\Term{morphism}{态射}，$E\in\operatorname{ob}D^b(X)$ 是\Term{cohomology}{上同调}\Term{quasi-coherent}{拟凝聚}且相对于 $f$ 具有\Term{finite}{有限}\Term{flat-amplitude}{平坦振幅}的\Term{complex}{复形}，$p\in\mathbf Z$。考虑以下条件：
\begin{enumerate}
\item[a)] \Term{functor}{函子} $M\mapsto R^pf_*(E\mathbin{\overset{\mathbf L}{\otimes}}_YM)$ 从 $\operatorname{Qcoh}(Y)$\footnote{$\operatorname{Qcoh}(Y)$ 是\Term{quasi-coherent}{拟凝聚} $\mathcal O_Y$-\Term{module}{模}的\Term{category}{范畴}。}到 $\operatorname{Qcoh}(Y)$\footnote{这是由 3.7 保证的。}为\Term{right-exact}{右正合}。
\item[a')] 对一切 $M\in\operatorname{ob}\operatorname{Qcoh}(Y)$，\Term{canonical}{典范}\Term{arrow}{箭头}
\[
R^pf_*(E)\otimes M\longrightarrow R^pf_*(E\mathbin{\overset{\mathbf L}{\otimes}}_YM)
\]
为\Term{isomorphism}{同构}。这里该\Term{arrow}{箭头}为\Term{canonical}{典范}\Term{arrow}{箭头}
\[
R^pf_*(E)\otimes M\longrightarrow H^p(\mathbf Rf_*(E)\mathbin{\overset{\mathbf L}{\otimes}}M)
\]
与 3.7 中的\Term{isomorphism}{同构}
\[
H^p(\mathbf Rf_*(E)\mathbin{\overset{\mathbf L}{\otimes}}M)\xrightarrow{\sim}R^pf_*(E\mathbin{\overset{\mathbf L}{\otimes}}_YM)
\]
的复合。
\item[b)] 从 $\operatorname{Qcoh}(Y)$ 到 $\operatorname{Qcoh}(Y)$ 的\Term{functor}{函子} $M\mapsto R^{p+1}f_*(E\mathbin{\overset{\mathbf L}{\otimes}}_YM)$ \Term{left-exact}{左正合}。
\item[b')] 对一切 $M\in\operatorname{ob}\operatorname{Qcoh}(Y)$，\Term{canonical}{典范}\Term{arrow}{箭头}
\[
R^{p+1}f_*(E\mathbin{\overset{\mathbf L}{\otimes}}_YM)\longrightarrow
\mathcal Hom(H^{-p-1}((\mathbf Rf_*(E))^\vee),M)
\]
为\Term{isomorphism}{同构}。这里该\Term{arrow}{箭头}为 3.7 中的\Term{isomorphism}{同构}
\[
R^{p+1}f_*(E\mathbin{\overset{\mathbf L}{\otimes}}_YM)\xrightarrow{\sim}
H^{p+1}(\mathbf Rf_*(E)\mathbin{\overset{\mathbf L}{\otimes}}M)
\]
与 5.1 $(*)$ 中的\Term{canonical}{典范}\Term{arrow}{箭头}
\[
H^{p+1}(\mathbf Rf_*(E)\mathbin{\overset{\mathbf L}{\otimes}}M)\longrightarrow
\mathcal Hom(H^{-p-1}((\mathbf Rf_*(E))^\vee),M)
\]
的复合。
\SourcePage{274}
\item[c)] 对任意\Term{base-change}{基变换} $Y'\to Y$，\Term{canonical}{典范}\Term{arrow}{箭头}
\[
R^pf_*(E)\otimes_Y\mathcal O_{Y'}\longrightarrow
H^p(\mathbf Rf_*(E)\mathbin{\overset{\mathbf L}{\otimes}}_Y\mathcal O_{Y'})
\]
为\Term{isomorphism}{同构}。
\end{enumerate}
则有 $\text{a)}\Longleftrightarrow\text{a')}\Longleftrightarrow\text{b)}\Longleftrightarrow\text{c)}\Longleftarrow\text{b')}$。
此外，若 $f$ \Term{proper}{固有}，$E$ 相对于 $f$ \Term{pseudo-coherent}{伪凝聚}（因而相对于 $f$ 为\Term{perfect-complex}{完美复形}），且 2.1 对 $f$ 成立（所以特别地，若 $f$ \Term{projective}{射影}或 $Y$ \Term{noetherian}{诺特}，此条件成立），则上述条件\Term{equivalence}{等价}。
\end{SGAresult}
\begin{proof}
由 3.7、5.1 和 5.2 (iv)，蕴含 $\text{a)}\Longleftrightarrow\text{a')}\Longleftrightarrow\text{b)}\Longleftarrow\text{b')}$ 立即得出；对于 $f$ \Term{proper}{固有}、$E$ 相对于 $f$ \Term{pseudo-coherent}{伪凝聚}的情形，由 2.1，b')$\Longleftrightarrow$b) 也同样得出。只须证明 a)$\Longleftrightarrow$c)。这由下述\Term{lemma}{引理}的第一个断言得出。
\begin{SGAresult}{\Term{lemma}{引理}}{5.3.1}
设 $X$ 是\Term{quasi-separated}{拟分离}\Term{scheme}{概形}，$E\in\operatorname{ob}D^-(X)$ 是\Term{cohomology}{上同调}\Term{local-scheme}{局部}\Term{bounded}{有界}且\Term{quasi-coherent}{拟凝聚}的\Term{complex}{复形}。则条件 5.1 a) \Term{equivalence}{等价}于：

c) 对 $X$ 上的每个\Term{scheme}{概形} $X'$，\Term{canonical}{典范}\Term{arrow}{箭头}
\[
H^p(E)\otimes_X\mathcal O_{X'}\longrightarrow H^p(E\mathbin{\overset{\mathbf L}{\otimes}}_X\mathcal O_{X'})
\]
为\Term{isomorphism}{同构}。

此外，若 $E$ \Term{pseudo-coherent}{伪凝聚}，则条件 5.1 a) 蕴含：对每个\Term{morphism}{态射} $X'\to X$，\Term{canonical}{典范}\Term{arrow}{箭头}
\[
H^{-p-1}(E^\vee)\otimes_X\mathcal O_{X'}\longrightarrow
H^{-p-1}((E\mathbin{\overset{\mathbf L}{\otimes}}_X\mathcal O_{X'})^\vee)\tag{5.3.1.1}
\]
\SourcePage{275}
为\Term{isomorphism}{同构}；此\Term{arrow}{箭头}是
\[
H^{-p-1}(E^\vee)\otimes_X\mathcal O_{X'}\longrightarrow
H^{-p-1}(E^\vee\mathbin{\overset{\mathbf L}{\otimes}}_X\mathcal O_{X'})
\]
与\Term{canonical}{典范}\Term{arrow}{箭头}（第 I 讲 7.1.2）
\[
H^{-p-1}(E^\vee\mathbin{\overset{\mathbf L}{\otimes}}_X\mathcal O_{X'})\longrightarrow
H^{-p-1}((E\mathbin{\overset{\mathbf L}{\otimes}}_X\mathcal O_{X'})^\vee)
\]
的复合。
\end{SGAresult}
\textbf{证明。} a)$\Rightarrow$c)。问题在 $X'$ 上是\Term{local-scheme}{局部}的，可设 $X$（相应地 $X'$）\Term{affine}{仿射}，其环为 $A$（相应地 $A'$）。又由第 II 讲 3.5，可设 $E=\widetilde F$，其中 $F$ 为 $A$-\Term{module}{模}的\Term{complex}{复形}。由假设（条件 a')），\Term{arrow}{箭头}
\[
H^p(F)\otimes_AA'\longrightarrow H^p(F\mathbin{\overset{\mathbf L}{\otimes}}_AA')
\]
为\Term{isomorphism}{同构}；因 $H^p(E)\otimes_X\mathcal O_{X'}$（相应地 $H^p(E\mathbin{\overset{\mathbf L}{\otimes}}_X\mathcal O_{X'})$）\Term{canonical}{典范}地等同于 $A'$-\Term{module}{模} $H^p(F)\otimes_AA'$（相应地 $H^p(F\mathbin{\overset{\mathbf L}{\otimes}}_AA')$）在 $X'$ 上定义的\Term{sheaf}{层}，结论即得。

c)$\Rightarrow$a)。由 5.2 (iv)，只须证明 c)$\Rightarrow$a'$_{\mathrm{qcoh}}$)。问题是\Term{local-scheme}{局部}的（因 $X$ \Term{quasi-separated}{拟分离}），可设 $X$ 为环 $A$ 的\Term{affine-scheme}{仿射概形}，$E=\widetilde F$，其中 $F$ 是 $A$-\Term{module}{模}的\Term{complex}{复形}。设 $M=\widetilde N$ 为\Term{quasi-coherent}{拟凝聚} $\mathcal O_X$-\Term{module}{模}。须证\Term{arrow}{箭头}
\[
H^p(F)\otimes_AN\longrightarrow H^p(F\mathbin{\overset{\mathbf L}{\otimes}}_AN)
\]
为\Term{isomorphism}{同构}。但取 $A'=D_A(N)$（EGA $0_{\mathrm{IV}}$ 18.2.3），$X'=\operatorname{Spec}(A')$，则对\Term{canonical}{典范}\Term{morphism}{态射} $X'\to X$ 应用 c)，即得此结论。

证明最后一个断言。问题在 $X'$ 上是\Term{local-scheme}{局部}的，由第 II 讲 2.2.10 b)，可设 $E$ 严格\Term{pseudo-coherent}{伪凝聚}。又因断言只涉及 $E$ 中\Term{length}{长度}\Term{bounded}{有界}的一段，所以经\Term{truncation}{截断}，可设 $E$ \Term{strictly-perfect}{严格完美}。此时 $(E\mathbin{\overset{\mathbf L}{\otimes}}_X\mathcal O_{X'})^\vee$ \Term{canonical}{典范}地等同于 $E^\vee\mathbin{\overset{\mathbf L}{\otimes}}_X\mathcal O_{X'}$（第 I 讲 7.1.2），\Term{arrow}{箭头} 5.3.1.1 可重写为\Term{canonical}{典范}\Term{arrow}{箭头}
\[
H^{-p-1}(E^\vee)\otimes_X\mathcal O_{X'}\longrightarrow
H^{-p-1}(E^\vee\mathbin{\overset{\mathbf L}{\otimes}}_X\mathcal O_{X'}).\qquad(*)
\]
\SourcePage{276}
若 $E$ 对整数 $p$ 满足条件 5.1 a)，则由 5.2 (iii)，$E^\vee$ 对整数 $-p-1$ 满足同一条件；由 5.3.1 的第一个断言，\Term{arrow}{箭头} $(*)$ 是\Term{isomorphism}{同构}。这就完成了 5.3.1 的证明，从而也完成了 5.3 的证明。
\end{proof}
\textbf{注记 5.3.2。} 特别地，当 $E$ 为\Term{bounded}{有界}\Term{complex}{复形}，各\Term{set-component}{分量}是相对于 $f$ \Term{flat}{平坦}且\Term{finitely-presented}{有限表示}的\Term{quasi-coherent}{拟凝聚} $\mathcal O_X$-\Term{module}{模}时，命题 5.3 适用；事实上，此时由 4.3.2，$E$ 相对于 $f$ 具有\Term{finite}{有限}\Term{perfect-amplitude}{完美振幅}。

\SourcePage{277}
\begin{SGAresult}{定义}{5.4}
设 $X$ 为\Term{topos}{拓扑斯}，记 $\operatorname{Pt}(X)$ 为它的点集（SGA 4 第 IV 讲 6.1）。设 $M$ 为\Term{ordered-set}{有序集}，$f:\operatorname{Pt}(X)\to M$ 为映射，$x\in\operatorname{Pt}(X)$。如果存在 $X$ 的一个\Term{pointed-object}{带 $x$ 点的对象} $U$，使 $U$ 的每个点 $y$ 均满足 $f(y)\leq f(x)$（相应地 $f(y)\geq f(x)$、$f(y)=f(x)$），就称 $f$ 在 $x$ 处\Term{upper-semicontinuous}{上半连续}（相应地\Term{lower-semicontinuous}{下半连续}、\Term{locally-constant}{局部常值}）。如果 $f$ 在每个点都如此，就称 $f$ \Term{upper-semicontinuous}{上半连续}（相应地\Term{lower-semicontinuous}{下半连续}、\Term{locally-constant}{局部常值}）。\footnote{如果 $X$ 来自一个\Term{topological-space}{拓扑空间}，这里就得到通常的概念。}
\end{SGAresult}
\textbf{例 5.4.1。} a) 记 $M_X$ 为 $X$ 上取值于 $M$ 的\Term{constant-sheaf}{常值层}，设 $s\in H^0(X,M)$。从 $\operatorname{Pt}(X)$ 到 $M$ 的映射 $x\mapsto s(x)=x^*(s)$ \Term{locally-constant}{局部常值}。特别地，假设 $X$ \Term{locally-ringed}{局部赋环}，$E$ 为 $X$ 上的\Term{perfect-complex}{完美复形}；则 $E$ 的\Term{rank}{秩}（第 I 讲 6.11）定义了 $\operatorname{Pt}(X)$ 上一个取值于 $\mathbf Z$ 的\Term{locally-constant}{局部常值}函数。

b) 假设 $X$ \Term{locally-ringed}{局部赋环}，$P$ 为\Term{finite-type}{有限型} $\mathcal O_X$-\Term{module}{模}。对 $x\in\operatorname{Pt}(X)$，记 $\mathcal O_{X,x}$ 为 $\mathcal O_X$ 在 $x$ 处的\Term{stalk}{茎}\Term{local-ring}{局部环}，$k(x)$ 为其\Term{residue-field}{剩余域}。则 $P_x\otimes_{\mathcal O_{X,x}}k(x)$ 在 $k(x)$ 上有\Term{finite}{有限}\Term{rank}{秩}，且函数
\[
x\longmapsto\operatorname{rg}_{k(x)}(P_x\otimes_{\mathcal O_{X,x}}k(x))
\]
\Term{upper-semicontinuous}{上半连续}。事实上，若 $P_x\otimes_{\mathcal O_{X,x}}k(x)$ 的\Term{rank}{秩}为 $n$，则由 \Term{nakayama-lemma}{Nakayama 引理}，存在\Term{epimorphism}{满态射} $\mathcal O_{X,x}^n\to P_x$；又因 $P$ \Term{finite-type}{有限型}，存在 $X$ 的\Term{pointed-object}{带 $x$ 点的对象} $U$ 及\Term{epimorphism}{满态射} $\mathcal O_U^n\to P|_U$，断言即得。

\begin{SGAresult}{\Term{lemma}{引理}}{5.5}
设 $(X,\mathcal O_X)$ 为\Term{locally-ringed-topos}{局部赋环拓扑斯}，$E\in\operatorname{ob}D^-(X)$，$p\in\mathbf Z$。沿用 5.4.1 b) 的记号，则：
\SourcePage{278}
\begin{enumerate}
\item[(i)] 若 $E$ 为\Term{perfect-complex}{完美复形}，则对 $x\in\operatorname{Pt}(X)$，各 $H^i(E_x\mathbin{\overset{\mathbf L}{\otimes}}_{\mathcal O_{X,x}}k(x))$ 是 $k(x)$ 上的\Term{finite}{有限}\Term{rank}{秩}\Term{vector-space}{向量空间}，仅\Term{finite}{有限}多个非零；函数
\[
x\longmapsto\sum_i(-1)^i\operatorname{rg}_{k(x)}H^i(E_x\mathbin{\overset{\mathbf L}{\otimes}}_{\mathcal O_{X,x}}k(x))
\]
在 $\operatorname{Pt}(X)$ 上\Term{locally-constant}{局部常值}。
\item[(ii)] 若 $E$ 是 $(p-1)$-\Term{pseudo-coherent}{伪凝聚}的，则对 $x\in\operatorname{Pt}(X)$，$H^p(E_x\mathbin{\overset{\mathbf L}{\otimes}}_{\mathcal O_{X,x}}k(x))$ 为 $k(x)$ 上的\Term{finite}{有限}\Term{rank}{秩}\Term{vector-space}{向量空间}，且函数
\[
x\longmapsto\operatorname{rg}_{k(x)}H^p(E_x\mathbin{\overset{\mathbf L}{\otimes}}_{\mathcal O_{X,x}}k(x))
\]
在 $\operatorname{Pt}(X)$ 上\Term{upper-semicontinuous}{上半连续}。若进一步从 $\operatorname{Mod}(X)$ 到 $\operatorname{Mod}(X)$ 的\Term{functor}{函子} $M\mapsto H^p(E\mathbin{\overset{\mathbf L}{\otimes}}M)$ \Term{exact}{正合}，则 $H^p(E)$ \Term{locally-free-finite-type}{有限型局部自由}，且上述函数\Term{locally-constant}{局部常值}。
\item[(iii)] 假设 $E$ 为 $p$-\Term{pseudo-coherent}{伪凝聚}（相应地 $(p-1)$-\Term{pseudo-coherent}{伪凝聚}）的，$x\in\operatorname{Pt}(X)$。若\Term{functor}{函子} $M\mapsto H^p(E_x\mathbin{\overset{\mathbf L}{\otimes}}_{\mathcal O_{X,x}}M)$ 在 $\mathcal O_{X,x}$-\Term{module}{模}的\Term{category}{范畴} $\operatorname{Mod}(x)$ 上\Term{right-exact}{右正合}（相应地\Term{left-exact}{左正合}），则存在 $X$ 的一个\Term{pointed-open}{带 $x$ 点的开对象} $U$，使\Term{functor}{函子} $M\mapsto H^p(E|_U\mathbin{\overset{\mathbf L}{\otimes}}M)$ 在 $\operatorname{Mod}(U)$ 上\Term{right-exact}{右正合}（相应地\Term{left-exact}{左正合}）。
\end{enumerate}
\end{SGAresult}
\begin{proof}
(i) 由第 I 讲 4.19.1，$E_x\mathbin{\overset{\mathbf L}{\otimes}}_{\mathcal O_{X,x}}k(x)$ 为 $k(x)$-\Term{module}{模}的\Term{perfect-complex}{完美复形}，故由第 I 讲 5.9 和 3.5，各 $H^i(E_x\mathbin{\overset{\mathbf L}{\otimes}}_{\mathcal O_{X,x}}k(x))$ 在 $k(x)$ 上有\Term{finite}{有限}\Term{rank}{秩}，且仅\Term{finite}{有限}多个非零。有
\[
\operatorname{rg}(E_x\mathbin{\overset{\mathbf L}{\otimes}}_{\mathcal O_{X,x}}k(x))
=\sum_i(-1)^i\operatorname{rg}_{k(x)}H^i(E_x\mathbin{\overset{\mathbf L}{\otimes}}_{\mathcal O_{X,x}}k(x)),\tag{5.5.1}
\]
\footnote{$\operatorname{rg}(\,\cdot\,)$ 表示\Term{perfect-complex}{完美复形}的\Term{rank}{秩}（第 I 讲 6.11）。}利用\Term{rank}{秩}函数的可\Term{additive}{加性}，逐步\Term{reduced}{约化}即可立即看出此式。结合 5.4.1 a)，断言 (i) 因而由
\SourcePage{279}
\Term{rank}{秩}与\Term{inverse-image}{逆像}的相容性（第 I 讲 6.12）得出。

(ii) 对第一个断言，经\Term{truncation}{截断}及\Term{translation}{平移}，可设 $E$ \Term{strictly-perfect}{严格完美}，$p=0$，且当 $i\notin[-1,+1]$ 时 $E^i=0$。对 $x\in\operatorname{Pt}(X)$，置 $E_x\mathbin{\overset{\mathbf L}{\otimes}}_{\mathcal O_{X,x}}k(x)=E_{k(x)}$。由 5.5.1，
\[
\operatorname{rg}_{k(x)}H^0(E_{k(x)})
=\operatorname{rg}_{k(x)}H^{-1}(E_{k(x)})
+\operatorname{rg}_{k(x)}H^1(E_{k(x)})+\operatorname{rg}(E_x),
\]
所以只须证明函数 $x\mapsto\operatorname{rg}_{k(x)}H^{-1}(E_{k(x)})$（相应地 $\operatorname{rg}_{k(x)}H^1(E_{k(x)})$）\Term{upper-semicontinuous}{上半连续}，因为函数 $x\mapsto\operatorname{rg}(E_x)=\operatorname{rg}(E_{k(x)})$ \Term{locally-constant}{局部常值}。对于函数 $x\mapsto\operatorname{rg}_{k(x)}H^1(E_{k(x)})$，这是 5.4.1 b) 的直接推论。\footnote{这里用到 $H^1(E_{k(x)})\simeq H^1(E)_{k(x)}$。}对于函数 $x\mapsto\operatorname{rg}_{k(x)}H^{-1}(E_{k(x)})$，可设 $E^1=0$，即以 $E'=(0\to E^{-1}\to E^0\to0)$ 代替 $E$。已知函数 $x\mapsto\operatorname{rg}_{k(x)}H^0(E'_{k(x)})$ \Term{upper-semicontinuous}{上半连续}，因函数
\[
x\longmapsto\operatorname{rg}_{k(x)}H^0(E'_{k(x)})-
\operatorname{rg}_{k(x)}H^{-1}(E'_{k(x)})=\operatorname{rg}(E'_x)
\]
\Term{locally-constant}{局部常值}，函数
\[
x\longmapsto\operatorname{rg}_{k(x)}H^{-1}(E'_{k(x)})=
\operatorname{rg}_{k(x)}H^{-1}(E_{k(x)})
\]
也\Term{upper-semicontinuous}{上半连续}。

现证明第二个断言。经\Term{truncation}{截断}，可设 $E$ \Term{strictly-perfect}{严格完美}，且 $i<p-1$ 时 $E^i=0$。于是 $\mathcal O_X$-\Term{module}{模} $Z'^i(E)=\operatorname{Coker}(E^{i-1}\to E^i)$ \Term{finitely-presented}{有限表示}。有\Term{exact-sequence}{正合列}
\[
0\longrightarrow H^p(E)\longrightarrow Z'^p(E)\longrightarrow E^{p+1}\longrightarrow Z'^{p+1}(E)\longrightarrow0.\qquad(*)
\]
由于\Term{functor}{函子} $M\mapsto H^p(E\mathbin{\overset{\mathbf L}{\otimes}}M)$ 既\Term{right-exact}{右正合}又\Term{left-exact}{左正合}，由 5.2 (iii)，$Z'^{p+1}(E)$ 和 $Z'^p(E)$ \Term{flat}{平坦}，故由第 I 讲 5.8.3，它们\Term{locally-free-finite-type}{有限型局部自由}。$(*)$ 的\Term{exact}{正合}性随即蕴含 $H^p(E)$ \Term{locally-free-finite-type}{有限型局部自由}，因而函数
\[
x\longmapsto\operatorname{rg}_{k(x)}H^p(E_{k(x)})=\operatorname{rg}(H^p(E))_x
\]
\SourcePage{280}
\Term{locally-constant}{局部常值}。这完成了 (ii) 的证明。

(iii) 可设 $E$ \Term{strictly-perfect}{严格完美}，且 $i<p$（相应地 $i<p-1$）时 $E^i=0$。由 5.2 (iii)，\Term{functor}{函子} $M\mapsto H^p(E_x\mathbin{\overset{\mathbf L}{\otimes}}_{\mathcal O_{X,x}}M)$ 在 $\operatorname{Mod}(x)$ 上\Term{right-exact}{右正合}（相应地\Term{left-exact}{左正合}），\Term{equivalence}{等价}于 $Z'^{p+1}(E_x)$（相应地 $Z'^p(E_x)$；沿用 (ii) 的记号）\Term{flat}{平坦}，因而\Term{free-finite-type}{自由且有限型}（因为它\Term{finitely-presented}{有限表示}）。因
\[
Z'^{p+1}(E_x)\simeq Z'^{p+1}(E)_x
\quad\text{（相应地 }Z'^p(E_x)\simeq Z'^p(E)_x\text{），}
\]
通常的论证（参见 EGA $0_{\mathrm I}$ 5.2.7）表明，存在 $X$ 的一个\Term{pointed-object}{带 $x$ 点的对象} $U$，使 $Z'^{p+1}(E|_U)$（相应地 $Z'^p(E|_U)$）\Term{free-finite-type}{自由且有限型}，结论即得。
\end{proof}
\begin{SGAresult}{推论}{5.6}
（EGA III 7.6.9。）设 $X$ 为\Term{scheme}{概形}，$p\in\mathbf Z$，$E\in\operatorname{ob}D^-(X)$ 为 $(p-1)$-\Term{pseudo-coherent}{伪凝聚}\Term{complex}{复形}。则：
\begin{enumerate}
\item[(i)] 对 $x\in X$，$H^p(E\mathbin{\overset{\mathbf L}{\otimes}}_Xk(x))$ 在 $k(x)$ 上有\Term{finite}{有限}\Term{rank}{秩}，且函数
\[
x\longmapsto\operatorname{rg}_{k(x)}H^p(E\mathbin{\overset{\mathbf L}{\otimes}}_Xk(x))
\]
在 $X$ 上\Term{constructible}{可构造}且\Term{upper-semicontinuous}{上半连续}。
\item[(ii)] 假设 $X$ \Term{quasi-separated}{拟分离}，$E$ 的\Term{cohomology}{上同调}\Term{quasi-coherent}{拟凝聚}且\Term{local-scheme}{局部}\Term{bounded}{有界}。若\Term{functor}{函子} $M\mapsto H^p(E\mathbin{\overset{\mathbf L}{\otimes}}M)$ 在\Term{quasi-coherent}{拟凝聚} $\mathcal O_X$-\Term{module}{模}的\Term{category}{范畴}上\Term{exact}{正合}，则函数 $x\mapsto\operatorname{rg}_{k(x)}H^p(E\mathbin{\overset{\mathbf L}{\otimes}}_Xk(x))$ \Term{locally-constant}{局部常值}，且 $H^p(E)$ \Term{locally-free-finite-type}{有限型局部自由}。反过来，若 $X$ \Term{locally-noetherian}{局部诺特}且\Term{reduced}{约化}，上述函数\Term{locally-constant}{局部常值}，则\Term{functor}{函子} $M\mapsto H^p(E\mathbin{\overset{\mathbf L}{\otimes}}M)$ 在\Term{quasi-coherent}{拟凝聚} $\mathcal O_X$-\Term{module}{模}的\Term{category}{范畴}上\Term{exact}{正合}。
\end{enumerate}
\end{SGAresult}
\begin{proof}
若 $X$ \Term{integral}{整}，\Term{generic-point}{泛点}为 $x$，则\Term{functor}{函子} $M\mapsto H^p(E_x\mathbin{\overset{\mathbf L}{\otimes}}_{\mathcal O_{X,x}}M)$ 在 $\mathcal O_{X,x}$-\Term{module}{模}的\Term{category}{范畴}上\Term{exact}{正合}（因为 $\mathcal O_{X,x}=k(x)$）；所以由 5.5 (iii)，存在 $x$ 的一个\Term{open-neighborhood}{开邻域} $U$，使\Term{functor}{函子} $M\mapsto H^p(E|_U\mathbin{\overset{\mathbf L}{\otimes}}M)$ 在 $\operatorname{Mod}(U)$ 上\Term{exact}{正合}，并且由 5.5 (ii)，函数
\SourcePage{281}
$y\mapsto\operatorname{rg}_{k(x)}H^p(E\mathbin{\overset{\mathbf L}{\otimes}}_Xk(y))$ 在 $U$ 上\Term{locally-constant}{局部常值}。由此立即得出 (i) 中的\Term{constructible}{可构造}性断言。其余断言由 5.5 (ii)、(iii) 和 5.2 (iv) 得出，只有 (ii) 的逆命题例外；对此，在\Term{localization}{局部化}和\Term{truncation}{截断}后归结到 $E$ \Term{strictly-perfect}{严格完美}的情形，再对\Term{functor}{函子} $T_{-p}=H^p(E\otimes\,\cdot\,)$ 应用 EGA III 7.6.2 和 7.6.5 即可。
\end{proof}
\begin{SGAresult}{命题}{5.7}
（EGA III 7.7.5 和 7.9.4。）设 $f:X\to Y$ 是\Term{quasi-compact}{拟紧}\Term{scheme}{概形}之间的\Term{proper}{固有}\Term{morphism}{态射}，$E\in\operatorname{ob}D^b(f)_{\underline{\mathrm{parf}}}$（4.2），$p\in\mathbf Z$。假定 2.1 成立，例如 $f$ \Term{projective}{射影}或 $Y$ \Term{noetherian}{诺特}，则：
\begin{enumerate}
\item[(i)] 对 $y\in Y$，各 $H^i(\mathbf Rf_*(E)\mathbin{\overset{\mathbf L}{\otimes}}_Yk(y))$ 为 $k(y)$ 上的\Term{finite}{有限}\Term{rank}{秩}\Term{vector-space}{向量空间}，仅\Term{finite}{有限}多个非零；函数
\[
y\longmapsto\sum_i(-1)^i\operatorname{rg}_{k(y)}H^i(\mathbf Rf_*(E)\mathbin{\overset{\mathbf L}{\otimes}}_Yk(y))
\]
在 $Y$ 上\Term{locally-constant}{局部常值}。
\item[(ii)] 函数
\[
y\longmapsto\operatorname{rg}_{k(y)}H^p(\mathbf Rf_*(E)\mathbin{\overset{\mathbf L}{\otimes}}_Yk(y))
\]
在 $Y$ 上\Term{constructible}{可构造}且\Term{upper-semicontinuous}{上半连续}。
\item[(iii)] 假设 $Y$ \Term{quasi-separated}{拟分离}。若\Term{functor}{函子} $M\mapsto R^pf_*(E\mathbin{\overset{\mathbf L}{\otimes}}_YM)$ 在\Term{quasi-coherent}{拟凝聚} $\mathcal O_Y$-\Term{module}{模}的\Term{category}{范畴}上\Term{exact}{正合}，则 $R^pf_*(E)$ 为\Term{locally-free-finite-type}{有限型局部自由} $\mathcal O_Y$-\Term{module}{模}，且 (ii) 中的函数\Term{locally-constant}{局部常值}。反之，若 $Y$ \Term{noetherian}{诺特}且\Term{reduced}{约化}，(ii) 中的函数\Term{locally-constant}{局部常值}，则\Term{functor}{函子} $M\mapsto R^pf_*(E\mathbin{\overset{\mathbf L}{\otimes}}_YM)$ 在\Term{quasi-coherent}{拟凝聚} $\mathcal O_Y$-\Term{module}{模}的\Term{category}{范畴}上\Term{exact}{正合}。
\end{enumerate}
\end{SGAresult}
\SourcePage{282}
\begin{proof}
由 4.8，$\mathbf Rf_*(E)$ 是\Term{perfect-complex}{完美复形}。因此，断言 (i)（相应地 (ii)）由 5.5 (i)（相应地 5.6 (i)）得出。最后，若 $M$ 为\Term{quasi-coherent}{拟凝聚} $\mathcal O_Y$-\Term{module}{模}，则由 3.7，\Term{canonical}{典范}\Term{arrow}{箭头}
\[
\mathbf Rf_*(E)\mathbin{\overset{\mathbf L}{\otimes}}_YM\longrightarrow
\mathbf Rf_*(E\mathbin{\overset{\mathbf L}{\otimes}}_YM)
\]
为\Term{isomorphism}{同构}，故断言 (iii) 由 5.6 (ii) 得出。
\end{proof}
\textbf{注记 5.9。}

a) 在 5.7 的情形，记 $i_y:X_y\to X$ 为 $X$ 在 $Y$ 的点 $y$ 上的\Term{fiber}{纤维}的\Term{embedding}{嵌入}。可以证明（参见 EGA III 6 和第 IV 讲 3.1.0），当 $f$ \Term{flat}{平坦}（相应地 $E$ \Term{bounded}{有界}且各\Term{set-component}{分量}\Term{quasi-coherent}{拟凝聚}、$Y$-\Term{flat}{平坦}）时，有 \Term{kunneth-formula}{Künneth 公式}
\[
H^i(\mathbf Rf_*(E)\mathbin{\overset{\mathbf L}{\otimes}}_Yk(y))
\xrightarrow{\sim}H^i(X_y,E_y),
\]
其中 $E_y=\mathbf Li_y^*(E)$（相应地 $i_y^*(E)$）。

\SourcePage{283}
b) 对\Term{complex-analytic-space}{复解析空间}也可建立类似本节的整套理论（\Term{commutativity}{交换性}质、\Term{semicontinuity}{半连续性}、\Term{continuity}{连续性}）。\footnote{而且这些性质正是首先在这一框架下提出的（例如见 Grauert [1]），随后才在 EGA III 7 中移植到\Term{algebraic-geometry}{代数几何}中。}其基本要素是：1）\Term{proper}{固有}\Term{morphism}{态射}的\Term{finiteness-theorem}{有限性定理}（“Grauert \Term{theorem}{定理}”）；2）\Term{projection-formula}{投影公式} 3.7 的类似物，它可表述如下：设 $f:X\to Y$ 为\Term{analytic-space}{解析空间}的\Term{morphism}{态射}，假定 $X$ 的\Term{dimension}{维数}有上界，设 $E\in\operatorname{ob}D^-(X)$、$M\in\operatorname{ob}D^-(Y)_{\mathrm{coh}}$（“coh”表示“\Term{cohomology}{上同调}\Term{coherent}{凝聚}”）。则\Term{canonical}{典范}\Term{arrow}{箭头}
\[
\mathbf Rf_*(E)\mathbin{\overset{\mathbf L}{\otimes}}M\longrightarrow
\mathbf Rf_*(E\mathbin{\overset{\mathbf L}{\otimes}}_YM)
\]
为\Term{isomorphism}{同构}（关于 $X$ 的假设保证 $f_*$ 的\Term{cohomological-dimension}{上同调维数}\Term{finite}{有限}，见 T.F II 5.12，故 $\mathbf Rf_*$ 在 $D(X)$ 上有定义）。

断言 2）与断言 1）不同，它本质上是显然的，可由类似 3.7 的证明中的逐步\Term{reduced}{约化}得出。又因 $\mathbf Rf_*$ 在 $D(X)$ 上有定义，论证更为自由；特别地，对 $E$ 无须任何假设（除了 $E\in\operatorname{ob}D^-(X)$）。

\textbf{参考文献}
\begin{description}
\item[{[1]}] Grauert, H.，\textit{Ein Theorem der Analytischen Garbentheorie\ldots}（解析层论中的一个\Term{theorem}{定理}……），Pub. IHES，第 5 号。
\item[{[H]}] Hartshorne, R.，\textit{Residues and Duality}（\Term{residue}{留数}与\Term{duality}{对偶}），Lecture Notes 第 20 号，Springer，1956。
\item[{[T.F]}] Godement, R.，\textit{Théorie des faisceaux}（层论），Hermann，1958。
\end{description}
